\SourceLeaf{400}%
\[
\begin{aligned}
\tfrac{\alpha y'+\beta(y''+y)-\eta}{\alpha x'+\beta(x''+x)-\xi}.\tfrac{y''-y}{x''-x}&=-1\\
\tfrac{\alpha y''+\beta(y+y')-\eta}{\alpha x''+\beta(x+x')-\xi}.\tfrac{y-y'}{x-x'}&=-1
\end{aligned}
\]
Hieraus folgt leicht durch Elimination
\[
\begin{aligned}
\xi&=\tfrac{(y-y')(y'-y'')(y''-y)(\alpha-\beta)}{y(x''-x')+y'(x-x'')+y''(x'-x)}\\
&\quad+\alpha.\tfrac{xy(x''-x')+x'y'(x-x'')+x''y''(x'-x)}{y(x''-x')+y'(x-x'')+y''(x'-x)}\\
&\quad+\beta.\tfrac{y(x''x''-x'x')+y'(xx-x''x'')+y''(x'x'-xx)}{y(x''-x')+y'(x-x'')+y''(x'-x)}
\end{aligned}
\]
\par
Der Werth von \(\eta\) folgt aus dem Werthe von \(\xi\), wenn man in diesem alle \(x,x',x''\) mit den entsprechenden \(y,y',y''\) vertauscht, wie man auch a priori leicht voraus sehen kann.

Die Coordinaten des Durchschnitts des ersten und letzten Perpendikels folgen, wie man leicht sieht, aus \(\xi\) und \(\eta\), wenn man \(x\) mit \(x'\), und \(y\) mit \(y'\) vertauscht, da aber dadurch \(\xi\) und \(\eta\) ihre Werthe nicht ändern, indem in beiden offenbar die Coordinaten der Punkte \(A,A',A''\) auf gleiche Art entriren, so ist klar, dass dieser zweite Durchschnittspunkt mit dem ersten zusammenfällt, und eben deshalb fällt der dritte Durchschnittspunkt mit den beiden ersten von selbst gleichfalls zusammen.

Für den Schwerpunkt ist übrigens offenbar
\[
n=\tfrac23
\]
also
\[
\alpha=\beta=\tfrac13
\]
und daher
\[
\begin{aligned}\xi&=\tfrac13(x+x'+x'')\\\eta&=\tfrac13(y+y'+y'')\end{aligned}
\]
\par
Für den Mittelpunkt des umschriebenen Kreises ist
\[
n=1
\]
also
\[
\alpha=0,\quad\beta=\tfrac12
\]
\par
Für den Durchschnittspunkt des Perpendikels aus \(A\) u.\,s.\,w. selbst ist
\[
n=0
\]
